\section{Quasi-Distributions vs. Pseudo Distributions, Similarities and Differences}

In a recent paper~\cite{Radyushkin:2017cyf}, an alternative interpretation has been
suggested for the Euclidean matrix elements proposed in LaMET, which provides
another method to extract the physical parton distribution from the same lattice calculations.
The starting point is the so-called Ioffe-time distribution ${\cal M}(\nu,\xi^2)$,
which is defined as the matrix elements of bilinear quark operators, separated with arbitrary
spacetime distances $\xi^\mu$, and in a hadron state with four-momentum $p^\mu$. As such, it
is an invariant quantity of $\nu= -p\cdot \xi$ (Ioffe-time), and space-time invariant $-\xi^2$.
In practice, such a Lorentz-invariant function can be obtained
on the Euclidean lattice
from the two-quark spatial correlation function,
\begin{equation}\label{eq:covitd}
h(zp_z,z^2,\Lambda) =\langle p|\overline{\psi}(z)\Gamma L(z,0)\psi(0)|p\rangle \ ,
\end{equation}
where $\Gamma$ is Dirac matrix, $L(z,0)$ the straight gauge link, and $\Lambda$
a renormalization scale or lattice cutoff scale. Thus, the matrix element
$h(zp_z,z^2,\Lambda)$ in the quasi distribution calculations in LaMET
can be viewed as a covariant Ioffe-time distribution $h(\nu,\xi^2,\Lambda)$ with $\nu = zp_z$, in space-like separation, $-\xi^2 = z^2$.

Radyushkin suggests to interpret the Ioffe-time distribution in a new frame, in which
$\xi^\mu = (\xi^+=0, \xi^-/\gamma, \xi_\perp)$, where we have used the light-cone
variable $\xi^\mu = (\xi^+,\xi^-,\xi_\perp)$,  and $ p^\mu = (p^+\gamma, 0, p^-=M^2/(p^+\gamma)$,
$M$ is the hadron mass, and $\gamma$ is a boost parameter.
With $\gamma\rightarrow\infty$, one recovers the infinite momentum frame.
In this new frame, the Ioffe time is $\nu = -p^+\xi^-$, and $-\xi^2= \xi_\perp^2$ is again space-like,
and they are $\gamma$ independent. And $h$ becomes a near light-cone correlation function,
\begin{equation}\label{eq:lccorrel}
h(p^+\xi^-, \xi_\perp, \Lambda) =\langle p|\overline{\psi}(0,\xi^-,\xi_\perp)\Gamma L(0,\xi^-,\xi_\perp;0)\psi(0)|p\rangle,
\end{equation}
where now the gauge link has both components along the light-cone and transverse
directions. Lorentz symmetry guarantees the quality of Eqs.~(\ref{eq:covitd}) and (\ref{eq:lccorrel}), with $\xi_\perp = z$,
and $\xi^- = -zp_z/p^+$. The Fourier transformation with respect to the Ioffe-time $\nu$ has the direct
interpretation as the light-cone momentum fraction $x$, therefore is bounded in $[-1,1]$
as always the case in the infinite momentum frame.

One can then introduce a ``pseudo''-distribution from the space-like correlation,
\begin{equation}
{\cal P}(y, \xi_\perp, \Lambda) = \frac{1}{2\pi} \int^\infty_{-\infty}d\nu\ e^{-iy\nu}\ h(zp_z,z^2, \Lambda)\ ,
\end{equation}
where $\xi_\perp^2 = z^2$. In this approach, the hadron momentum on lattice
is re-interpreted as the Ioffe time at a fixed $\xi_\perp$.
The physics of the pseudo-distribution is that it is a Fourier transformation
of a type of transverse-momentum dependent parton distributions. Since there
is a gauge link between the quark fields going along the transverse direction
in Eq.~(\ref{eq:lccorrel}), in $A^+=0$ guage, there is an infinite number of transversely polarized
gluons involved in the distributions. Therefore the parton picture of the pseudo-distribution
is not simple. In particular, the momentum fraction $y$ contains not only the contribution
from a single quark, but also that from transversely-polarized physical gluons. $y$ approaches
the quark light-cone fraction only in the limit of $\xi_\perp=0$.

In this alternative interpretation of the Euclidean matrix element,
the physical quark distribution is recovered by putting the quarks entirely on the light-cone,
i.e., when the transverse coordinate of the quark $\xi_\perp = z$ goes to zero.
In this limit, to get finite Ioffe-time $\nu = zp_z$, one again has to let
$p_z \rightarrow\infty$, the same as used in LaMET. Thus, the requirement to
get a precision parton distribution is exactly the same in the two approaches.

In practice, one has to work with small but nonvanishing $\xi_\perp$ to avoid light-cone divergences.
One can then establish a similar factorization theorem as in LaMET (``a small-distance effective theory''),
\begin{equation}
{\cal P}(y, \xi_\perp,\tilde{\mu}) = \int \frac{dx}{x}
C(y/x, \xi_\perp\tilde{\mu},\tilde{\mu}/\mu) q(x, \mu) + {\cal O}(\xi^2_\perp),
\label{fac2}
\end{equation}
where $q$ is a physical parton distribution, $C$ is a matching coefficient that depends
on the specific lattice regularization, $\tilde{\mu}$ and $\mu$ are the renormalization scales of the Ioffe-time distribution and physical PDF.
To have continuum limit, $\xi_\perp$ shall be much
larger than lattice spacing $a$, but much smaller than $1/\Lambda_{\text{QCD}}$. The corrections are in terms of small $\xi_\perp$ expansion.
It would be interesting to extract the parton distributions through the above
factorization formula at a fixed $\xi_\perp$.

To demonstrate the above factorization theorem, let us take the isovector unpolarized quark distribution as an example. In the Feynman gauge, we calculate the quark matrix elements of the quasi and physical PDFs at one-loop order in dimensional regularization with $D=4-2\epsilon$. The external quark state is chosen to be onshell and massless, so the UV and collinear divergences are regularized by $1/\epsilon_{UV}\ (\epsilon_{UV}>0)$ and $1/\epsilon_{IR}\ (\epsilon_{IR}<0)$ respectively. We compute the same Feynman diagrams in Ref.~\cite{Xiong:2013bka} in the coordinate space to obtain $h^{(1)}(zp^z,z^2,\mu,\epsilon)$ (for $\Gamma=\gamma^0$) and $q^{(1)}(p^+z^-,\mu,\epsilon)$, where $\mu$ can be regarded as the renormalization scale. The Feynman rules are not conventional as there is a Fourier transform of one loop momentum component in each of the irreducible diagrams, for instance,
\begin{align}
\mu^{2\epsilon}\int {d^d k\over (2\pi)^d} {1\over k^2(p-k)^2} e^{-i k^z z} &=  i\mu^{2\epsilon}\int_0^1 du \int {d^d k\over (2\pi)^d} {e^{-i k^z z}\over k_E^4} e^{ -i u p^z z}\ .
\end{align}
Using Schwinger parametrization, we can evaluate the above integral analytically,
\begin{align}
&i\mu^{2\epsilon}\int_0^1 du\int {d^d k\over (2\pi)^d} {e^{-i k^z z}\over k_E^4} e^{-i u p^z z}\nonumber\\
&= i\mu^{2\epsilon}\int_0^1 du\int_0^\infty d\alpha\ \alpha\int {d^d k\over (2\pi)^d} e^{-\alpha k_E^4-i k^z z} e^{-i u p^z z} \nonumber\\
&= {i\over 16\pi^2}(\pi z^2\mu^2)^{\epsilon_{IR}} \Gamma(-\epsilon_{IR})\int_0^1 du\  e^{-iu p^z z}\ .
\end{align}
The complete results for $h^{(1)}(zp^z,z^2,\mu,\epsilon)$ and $q^{(1)}(p^+z^-,\mu,\epsilon)$ are
\begin{align}
&h^{(1)}(zp^z,z^2,\mu,\epsilon)\nonumber\\
=&{\alpha_sC_F\over 2\pi}\int_0^1 du \left[\left(1-u+\left({2u\over 1-u}\right)_+\right)\left(-\ln(z^2 \mu^2) - {1\over \epsilon'_{IR}}\right)\right.\nonumber\\
&\ \ \ \ \ \ \ \ + (1-u)\Big]e^{-iu p^z z}\nonumber\\
& + {\alpha_sC_F\over 2\pi}(izp^z)\int_0^1 du \int_0^1dt\ (2-u) (-\ln t^2) e^{-i(1-tu) p^z z} \nonumber\\
&  +\left[(1+\delta Z_\psi)+ {\alpha_sC_F\over 2\pi}\left(2\ln(z^2\mu^2) + {2\over \epsilon'_{UV}}+2\right)\right] e^{-i p^z z} \ ,
\end{align}

\begin{align}
&q^{(1)}(p^+z^-,\mu,\epsilon)\nonumber\\
=&{\alpha_sC_F\over 2\pi}\int_0^1 du \left(1-u+\left({2u\over 1-u}\right)_+\right)\left({1\over \epsilon'_{UV}}- {1\over \epsilon'_{IR}}\right)e^{iu p^+ z^-}\nonumber\\
& + (1+\delta Z_\psi)e^{i p^+ z^-} \ ,
\end{align}
where $Z_\psi = 1+ \delta Z_\psi$ is the quark wavefunction renormalization constant, and we have made the substitutions $1/\epsilon'_{UV,IR} = 1/\epsilon_{UV,IR}+\gamma_E+\ln\pi$. To guarantee vector current conservation, $Z_\psi$ is given in the on-shell scheme,
\beq
Z_\psi = 1 - {\alpha_sC_F\over 2\pi}{1\over2}\left({1\over \epsilon'_{UV}}- {1\over \epsilon'_{IR}}\right) + O(\alpha_s^2)\ .
\eeq
Then we Fourier transform the Ioffe time $p^z z$ or $p^+z^-$ into $x$, and obtain
\begin{align} \label{eq:quasipdf}
&{\cal P}^{(1)}(x,\xi_\perp,\mu,\epsilon)\nonumber\\
=&{\alpha_sC_F\over 2\pi}\left[\left({1+x^2\over 1-x}\right)_+\left(-\ln(\xi_\perp^2 \mu^2) - {1\over \epsilon'_{IR}}-1\right)  - \left(4\ln(1-x)\over 1-x\right)_+  \right.\nonumber\\
&+2(1-x)\Big]\theta(x)\theta(1-x)\nonumber\\
&  +\left[ 1+ {\alpha_sC_F\over 2\pi}\left({3\over2}\ln(\xi_\perp^2\mu^2) + {3\over2}{1\over \epsilon'_{UV}}+{3\over2}\right)\right]\delta(1-x) \ ,
\end{align}
whereas the corresponding light-cone PDF is
\begin{align} \label{eq:lcpdf}
q^{(1)}(x,\mu,\epsilon) &= \delta(1-x) + {\alpha_s\over 2\pi}\left({1+x^2\over 1-x}\right)_+\left({1\over\epsilon'_{UV}} - {1\over \epsilon'_{IR}}\right)\ .
\end{align}
The support in the pseudo distribution ${\cal P}^{(1)}(x,\xi_\perp,\mu,\epsilon)$ is restricted to $0<x<1$, as expected.
If we renormalize both Eq.~(\ref{eq:quasipdf}) and Eq.~(\ref{eq:lcpdf}) in the $\overline{\text{MS}}$ scheme, then the matching coefficient in Eq.~(\ref{fac2}) is read off as
\begin{align}
&C\left(y,\xi_\perp\mu\right)\nonumber\\
=& \left[1+{\alpha_sC_F\over 2\pi}\left({3\over2}\ln(\xi_\perp^2\mu^2)+{3\over2}\right)\right] \delta(1-y)\nonumber\\
& + {\alpha_sC_F\over 2\pi} \left[-\left({1+y^2\over 1-y}\right)_+\left(\ln(\xi_\perp^2 \mu^2) +1 \right)- \left(4\ln(1-y)\over 1-y\right)_+ \right.\nonumber\\
&+2(1-y)\Big]\theta(y)\theta(1-y)\ .
\end{align}
\noindent
Similar factorization formula can be extended to the case of lattice regularization and non-perturbative renormalization.

Instead of using the factorization formula in Eq.~(\ref{fac2}) to extract parton distribution
from lattice matrix element calculated at a fixed $\xi_\perp$, Ref. \cite{Radyushkin:2017cyf}
focused on the possible approximate factorized property of the Ioffe-time distribution,
i.e. ${\cal M}(\nu, z) \sim  g(\nu)h(z)$. If so, $g(\nu) \sim {\cal M}(\nu, z)/h(z)$ will have approximate
scaling at different $z$. One can then get the small $z$
Ioffe-time distribution from a large $z$ calculation, which does not need a very large
momentum to produce a large Ioffe time. Indeed, an exploratory calculation on lattice
shows that this special factorization or scaling works quite well~\cite{Orginos:2017kos}.

While such a scaling or factorization behavior is interesting, just like parton-hadron duality
in deep-inelastic scattering~\cite{Ji:1994br} allowing the extraction of parton distributions at small $Q^2$, it is of limited use in actual precision calculations
of parton distributions. The scaling in the exploratory calculations
was observed in a limited kinematic domain with limited precision.
The parton distributions involve Ioffe time at all sizes, particularly small
$x$ distribution involves very large Ioffe time, where the scaling becomes much harder to test
and the present result is certainly incorrect. To get the correct
large Ioffe-time behavior, one needs calculation at small $\xi_\perp$ and larger nucleon momenta, the same
as required in LaMET. Moreover, the new $\xi_\perp$ dependence factorization must be corrected
for precision calculations. Since there is no systematic approach to correct for the violation
of this factorization, such a phenomenological approach
will not be a replacement for rigorous approaches as in Eq.~(\ref{fac2}), where
corrections can be quantified.

It of course will be interesting to compare the parton distributions from two seemingly
different factorization approaches (large momentum vs. small distance), when extracted from the same lattice matrix elements.
It is important to quantify the errors of extraction in both approaches consistently.
We emphasize that LaMET is a much more general framework to calculate parton physics
with simple physical pictures as provided by Feynman. It has a well-defined recipe on
how to systematically calculate all parton observables with quantifiable errors,
including multi-parton amplitudes and correlations.
